\chapter{Exercises with answers}\label{app:exercises}

Do them on paper first. Python is allowed for the last ones.

\begin{enumerate}
  \item Compute the odds and the OR for downvote shares 20\% and 10\%, then for 2\% and 1\%. What happens to the OR relative to the risk ratio?
  \item In Section~\ref{sec:h1}'s toy, change table Y to $(2,8,8,32)$ so its inside OR is 1.0 still, then to $(4,6,8,32)$. Predict whether MH goes up or down in each case and check with \file{mh\_toy.py}.
  \item With prior Beta$(1,2)$, a pair has 3 positive and 12 negative votes. Compute $E[p]$ and the edge sign. Then with 3 positive and 4 negative votes.
  \item A user has 2 replies, 1 hostile. The population share is 0.20 and $M=2$. What is the shrunk share? What if they had 20 replies, 10 hostile?
  \item KM: users leave at years 1, 3, 3 and are censored at 2 and 5. Compute $S(t)$ at each leaving time.
  \item Explain in two sentences why H1's interval is too narrow.
  \item Explain why the share of negative edges in H2 is not directly comparable between real and null pairs.
  \item A hazard ratio is 0.95 for ``got any reply''. State it in a sentence that a non-statistician understands.
  \item Why must the landmark Cox model exclude users who left before day 30?
  \item Name two ways the long-lived group can look ``calmer'' without calm users being more likely to stay.
\end{enumerate}

\section*{Answers}
\begin{enumerate}
  \item 20\%: odds 0.25; 10\%: 0.111; OR $=2.25$, risk ratio 2. 2\%: 0.0204; 1\%: 0.0101; OR $=2.02$, risk ratio 2. The OR approaches the risk ratio as the event gets rare, and exaggerates it when the event is common.
  \item $Y=(2,8,8,32)$ has inside OR $2\cdot32/(8\cdot8)=1$, like the original Y, but the downvote shares are higher, so Y's weight changes and MH falls from 1.58 to 1.44. $Y=(4,6,8,32)$ has inside OR $4\cdot32/(6\cdot8)=2.67$, above X's 2, so MH rises to 2.33. (Values computed with \file{mh\_toy.py}.)
  \item $3$ positive, $12$ negative: $E[p]=(3+1)/(15+3)=0.222$, negative. $3$ positive, $4$ negative: $E[p]=(3+1)/(7+3)=0.40$; not below 0.4, so no edge (and the pair has 7 votes, enough for the $\ge5$ rule).
  \item $(1+2\cdot0.2)/(2+2)=0.35$. With 10 of 20: $(10+0.4)/(22)=0.473$.
  \item At $t=1$: 5 at risk, 1 leaves, $S=0.8$. At $t=3$: users at risk are those with time $\ge3$: the two leaving at 3 and the one censored at 5, so 3 at risk, 2 leave, $S=0.8\cdot(1-2/3)=0.267$. The user censored at 2 left the risk set before $t=3$ without being counted as leaving.
  \item Many votes on different comments come from the same voters, and the interval treats comments as independent, so the effective sample size is overstated. The voter-level analysis gives a smaller effect with wider uncertainty.
  \item Real pairs more often reach five votes (repliers have a history with the people they answer), so the two shares come from different sets of pairs; a compositional difference could drive part of the gap.
  \item ``Among users who are still active at day 30, those who received a reply in their first month left at a slightly lower rate than those who did not (5\% lower, after accounting for cohort and other early exposures).''
  \item Because exposures are measured in the first 30 days; a user who left earlier would have an incomplete measurement, and including them would credit the exposure with time they could not have left.
  \item (a) The group is composed of survivors from earlier cohorts, which differ in tone; (b) quiet users who vote little receive few downvotes and also keep a long, low-intensity presence; (c) the share measures the \emph{rate} of downvotes which depends on how often the user writes controversial things, not on being calm.
\end{enumerate}
